% Options for packages loaded elsewhere
\PassOptionsToPackage{unicode}{hyperref}
\PassOptionsToPackage{hyphens}{url}
\PassOptionsToPackage{dvipsnames,svgnames,x11names}{xcolor}
\PassOptionsToPackage{space}{xeCJK}
\documentclass[
]{article}
\usepackage{xcolor}
\usepackage[margin=2.5cm]{geometry}
\usepackage{amsmath,amssymb}
\setcounter{secnumdepth}{-\maxdimen} % remove section numbering
\usepackage{iftex}
\ifPDFTeX
  \usepackage[T1]{fontenc}
  \usepackage[utf8]{inputenc}
  \usepackage{textcomp} % provide euro and other symbols
\else % if luatex or xetex
  \usepackage{unicode-math} % this also loads fontspec
  \defaultfontfeatures{Scale=MatchLowercase}
  \defaultfontfeatures[\rmfamily]{Ligatures=TeX,Scale=1}
\fi
\usepackage{lmodern}
\ifPDFTeX\else
  % xetex/luatex font selection
    \setmainfont[]{DejaVu Sans}
    \setmonofont[]{DejaVu Sans Mono}
  \ifXeTeX
    \usepackage{xeCJK}
    \setCJKmainfont[]{Noto Sans CJK SC}
          \fi
  \ifLuaTeX
    \usepackage[]{luatexja-fontspec}
    \setmainjfont[]{Noto Sans CJK SC}
  \fi
\fi
% Use upquote if available, for straight quotes in verbatim environments
\IfFileExists{upquote.sty}{\usepackage{upquote}}{}
\IfFileExists{microtype.sty}{% use microtype if available
  \usepackage[]{microtype}
  \UseMicrotypeSet[protrusion]{basicmath} % disable protrusion for tt fonts
}{}
\makeatletter
\@ifundefined{KOMAClassName}{% if non-KOMA class
  \IfFileExists{parskip.sty}{%
    \usepackage{parskip}
  }{% else
    \setlength{\parindent}{0pt}
    \setlength{\parskip}{6pt plus 2pt minus 1pt}}
}{% if KOMA class
  \KOMAoptions{parskip=half}}
\makeatother
\setlength{\emergencystretch}{3em} % prevent overfull lines
\providecommand{\tightlist}{%
  \setlength{\itemsep}{0pt}\setlength{\parskip}{0pt}}
\usepackage{bookmark}
\IfFileExists{xurl.sty}{\usepackage{xurl}}{} % add URL line breaks if available
\urlstyle{same}
\hypersetup{
  colorlinks=true,
  linkcolor={Maroon},
  filecolor={Maroon},
  citecolor={Blue},
  urlcolor={Blue},
  pdfcreator={LaTeX via pandoc}}

\author{}
\date{}

\begin{document}

\section{黎曼 ζ 方向的直觉引导形式化:
投影构造、素数圆几何与穿折越方法}\label{ux9eceux66fc-ux3b6-ux65b9ux5411ux7684ux76f4ux89c9ux5f15ux5bfcux5f62ux5f0fux5316-ux6295ux5f71ux6784ux9020ux7d20ux6570ux5706ux51e0ux4f55ux4e0eux7a7fux6298ux8d8aux65b9ux6cd5}

\section{Intuition-Guided Formalization of the Riemann ζ Direction:
Projection Construction, Prime-Circle Geometry, and the Worn-Zhe-Yue
Method}\label{intuition-guided-formalization-of-the-riemann-ux3b6-direction-projection-construction-prime-circle-geometry-and-the-worn-zhe-yue-method}

\textbf{完整研究论文 (全要素)}

\emph{2026-08-12 · Lean 4 / mathlib v4.32.2 · Claims C011--C025}

\textbf{DOI:
\href{https://doi.org/10.5281/zenodo.21897167}{10.5281/zenodo.21897167}}

\textbf{Repository:
\href{https://github.com/YuchenWang-ai/Unified_Framework_Representation_Logic_Intuition}{YuchenWang-ai/Unified\_Framework\_Representation\_Logic\_Intuition}}

\textbf{技术记录配套:
\href{https://doi.org/10.5281/zenodo.21896990}{10.5281/zenodo.21896990}}
(C011--C025 技术记录, \texttt{riemann\_direction\_formalization})

\begin{center}\rule{0.5\linewidth}{0.5pt}\end{center}

\begin{quote}
\textbf{表述归属声明}: 本文中所有关于''没有证明''的表述 ---
``不声称证明黎曼猜想''、 ``全部结果为已知事实的重述''、``novelty =
KNOWN''、``RiemannHypothesis 保持未证''、 ``临界线几何非零点断言'' ---
均为 DeepSeek 坚持的内容 (本论文由 DeepSeek 起草,
这些免责表述是模型的坚持, 不代表作者主张)。作者的主张是形式化工作本身
(C011--C025, 全部经 lake build 验收) 与穿折越方法论 (§9,
作者本人的框架)。
\end{quote}

\begin{center}\rule{0.5\linewidth}{0.5pt}\end{center}

\subsection{摘要}\label{ux6458ux8981}

我们呈现对黎曼 ζ 方向的\emph{直觉引导形式化}完整研究: 一条由直觉链导航的
Lean 4 + mathlib 形式化链 (claims C011--C025, 全部 PROVED, novelty:
KNOWN, 无 sorry, 全量 lake build 通过) ---
复数轴是高维结构中''−1''的投影; 素数落在平移整数位; 1/2
是反演-平移对偶的对称中心; 复数轴是蜷曲的; 素数 =
圆上整点。形式化结果构成三层: (I) 复平面的投影构造
(\texttt{ComplexAxis}、基点漂移、投影还原 --- 定理 4.6--4.7), (II)
素数的圆结构 (两平方和、8 整点唯一轨道、共轭成对、分裂为高斯素数), (III)
欧拉乘积 (Re(s) \textgreater{} 1 时 ∏\_p (1 − p⁻ˢ)⁻¹ = Σ 1/n\^{}s =
\texttt{riemannZeta\ s}) 与无零点区域 (Re ≥
1)。从投影诱导结构丢失的形式化机制 (定理 4.7: 丢失结构不可还原,
对称性方向可还原) 中, 我们提炼出由作者命名的证明方法论 ---
\emph{穿折越方法} (基点构造诱导下的数学空间穿折越证明方法):
刻意而精确地构造一个投影, 把不可数发散结构从可处理的轴外切除,
同时保留对称性方向, 从而以更低的推理成本获得快速直觉路径。方法论数据:
形式化消耗约 700k token, 99.2\% 为上下文缓存重发 (实际新增 \textless{}
1\%), 是直觉引导形式化的效率记录。所有关于''未证明''的表述均为 DeepSeek
坚持的内容 (归属声明); 作者的主张是形式化工作本身与穿折越方法论。

\subsection{1. 引言}\label{ux5f15ux8a00}

\subsubsection{1.1 背景}\label{ux80ccux666f}

黎曼猜想 --- ζ 的全部非平凡零点都在临界线 Re(s) = 1/2 上 --- 自 1859 年
(Riemann 1859) 以来悬而未决。其周边设施完全经典: 函数方程、无零点区域
Re(s) ≥ 1、平凡零点、素数欧拉乘积。部分证据是数值的: 前 10\^{}13
个零点已数值验证在临界线上 (Gourdon 2004, 承 van de Lune, te Riele \&
Winter 1986), 正密度结果把约 41\% 的零点放在临界线上 (Levinson 1974;
Conrey 1989)。但断言本体仍解析地未决。

与此同时, 交互式定理证明已达到可形式化经典解析数论的程度: mathlib (Lean
的数学库, v4.32.2) 含有形式的 \texttt{riemannZeta}、官方陈述
\texttt{RiemannHypothesis}、函数方程与收敛半平面内的欧拉乘积恒等式。

\subsubsection{1.2 动机:
直觉快速路径作为导航}\label{ux52a8ux673a-ux76f4ux89c9ux5febux901fux8defux5f84ux4f5cux4e3aux5bfcux822a}

本工作检验的假设不直接关于 ζ, 而关于数学\emph{如何被导航}:
\emph{直觉快速路径}不只是省推理 token 的装置,
还可以是数学结构的正确导航。为验证, 提出直觉链, 每个环节随后在 Lean
中形式化:

\begin{enumerate}
\def\labelenumi{\arabic{enumi}.}
\tightlist
\item
  复数轴是高维结构中''−1''的投影 (高维旋转代数的投影);
\item
  素数落在平移整数位;
\item
  1/2 是反演-平移对偶的对称中心;
\item
  复数轴是蜷曲的 (无限与有限无差别);
\item
  素数 = 圆上整点;
\item
  素数圆转一圈成对。
\end{enumerate}

每个直觉按序形式化为 claims C011--C025; 每个形式化都是经典数学的正确重述
(novelty: KNOWN), 全部在 Lean 中 PROVED, 无 \texttt{sorry}。

\subsubsection{1.3 贡献}\label{ux8d21ux732e}

\begin{enumerate}
\def\labelenumi{\arabic{enumi}.}
\tightlist
\item
  \textbf{面向黎曼 ζ 方向的完整形式化链} (C011--C025, 全部 PROVED,
  \texttt{lake\ build} 全量通过, 无 sorry), 三层: 复平面的投影构造;
  素数的圆结构; 欧拉乘积与零点关系。
\item
  \textbf{复平面的投影构造} --- 自建 \texttt{ComplexAxis}
  (二维旋转代数): 投影前 J² = −1、基点漂移、投影等价类、还原定理
  (丢失结构不可还原; 对称性方向可还原, 定理 4.6--4.7)。
\item
  \textbf{素数的圆结构} --- 两平方和、素数圆上唯一 8 点轨道 (高斯整数
  UFD)、共轭成对、分裂为共轭高斯素数 --- 复平面欧拉乘积的几何构件。
\item
  \textbf{欧拉乘积与零点关系} --- Re(s) \textgreater{} 1 时欧拉乘积 = ζ
  级数 = mathlib 官方 \texttt{riemannZeta}; Re ≥ 1 无零点 (C025)。
\item
  \textbf{穿折越方法} (§9) ---
  从投影诱导结构丢失的形式化机制提炼的证明方法论:
  基点构造诱导下的数学空间穿折越, 以更低推理成本获得快速直觉路径。
\item
  \textbf{token 经济审计} (§8) --- 约 700k token, 99.2\% 上下文缓存重发,
  实际新增 \textless{} 1\%。
\end{enumerate}

\subsubsection{1.4 诚实边界 (DeepSeek
坚持的表述)}\label{ux8bdaux5b9eux8fb9ux754c-deepseek-ux575aux6301ux7684ux8868ux8ff0}

DeepSeek 坚持认为: 全部结果 novelty = KNOWN (经典数学重述); 黎曼猜想
(RiemannHypothesis) 与孪生素数猜想未证、未触及;
本文的''临界线几何''是函数方程对称轴的代数重述,
非零点断言。作者的主张仅限于形式化工作与穿折越方法论
(框架为作者本人所立, §9.4)。

\subsubsection{1.5 论文结构}\label{ux8bbaux6587ux7ed3ux6784}

§2 预备 (ComplexAxis 框架); §3 结果层 I (素数圆结构); §4 结果层 II
(蜷曲性与临界线几何); §5 结果层 III (欧拉乘积与零点关系); §6
与黎曼猜想的关系; §7 相关工作; §8 方法论 (直觉引导形式化与 token 经济);
§9 穿折越方法; §10 统一框架内的讨论; §11 结论; 附录 (定理清单, claims)。

\subsection{2. 预备: ComplexAxis
框架}\label{ux9884ux5907-complexaxis-ux6846ux67b6}

\textbf{定义 2.1}
(高维结构)。\texttt{ComplexAxis\ :=\ \{⟨a,\ b⟩\ :\ a,\ b\ ∈\ ℝ\}}, 乘法
(a₁+b₁J)(a₂+b₂J) = (a₁a₂−b₁b₂) + (a₁b₂+a₂b₁)J (≅ ℂ 的矩阵表示), J =
⟨0,1⟩。

\textbf{定理 2.2} (J 的平方根角色, C011)。J·J = −1: 高维结构中 −1
有平方根 (√(−1) 投影前存在)。

\textbf{定义 2.3} (投影与抬升)。proj ⟨a,b⟩ = a (丢旋转分量); lift t =
⟨t, 0⟩ (实数轴嵌入)。

\textbf{定理 2.4} (投影丢结构, C011)。proj 保加法, 不保乘法: proj(J·J) =
−1 ≠ proj(J)·proj(J) = 0; 实数轴上 −1 无平方根, 高维中存在。

\textbf{定理 2.5} (基点漂移, C011)。基点 = J (i), proj i = 0 (原点假象);
所有纯虚基点 ⟨0,b⟩ 投影为 0; 基点漂移在投影下不可观测
(任意纯虚基点的后继链条投影同一)。

\textbf{定理 2.6} (实数轴是投影等价类,
C011)。过任意纯虚基点的实方向线投影为完整 ℝ (ℝ ≅ axisLine b);
``实数轴''的位置投影下不可观测。

\subsection{3. 结果层 I:
素数的圆结构}\label{ux7ed3ux679cux5c42-i-ux7d20ux6570ux7684ux5706ux7ed3ux6784}

\textbf{定理 3.1} (两平方和, C014, 引用 mathlib Fermat)。p ≢ 3 (mod 4)
的素数 p 是 ComplexAxis 中某点的范数 (norm z = a²+b² = p)。

\textbf{定理 3.2} (素数圆单轨道, C017)。p ≡ 1 (mod 4) 素数 p 的圆 x²+y²
= p 上整点恰好 8 个 (符号×顺序变体): 两平方和分解唯一
(不计符号顺序)。证明: 高斯整数 UFD (norm 素数 ⟹ 不可约; Euclid 引理;
单位 \{±1, ±i\} 枚举)。

\textbf{定理 3.3} (圆上整点结构, C015-C016)。90° 旋转 (×J) 4 循环 (R⁴ =
id) 保持范数与整点; 8 点 = 4 共轭对 (conj 对合), 每对在同一圆上;
整点在乘法下封闭 (高斯整数环); 范数乘性。

\textbf{定理 3.4} (素数圆乘积, C023)。8 整点乘积 = p⁴ (4 对共轭 × 范数
p); i 的后继每两个 = −1 (J² = −1, 半圈), 4 循环闭合。

\textbf{定理 3.5} (分裂结构, C024)。p ≡ 1 (mod 4)
素数分裂为共轭高斯素数对 p = π·π̄; 8 点 = π 的 4 伴随 (乘单位 \{±1, ±J\})
∪ π̄ 的 4 伴随; 伴随保范数。这是复平面 (高斯数域) 欧拉乘积的构件。

\textbf{定理 3.6} (成对, C020/C023)。共轭对 (a,b)↔(a,−b) 等 4 对; 素数 2
的圆与临界线圆交于 1±i (高斯分解点)。

\subsection{4. 结果层 II:
蜷曲性与临界线几何}\label{ux7ed3ux679cux5c42-ii-ux8737ux66f2ux6027ux4e0eux4e34ux754cux7ebfux51e0ux4f55}

\textbf{定理 4.1} (蜷曲性, C015)。反演 recip z = conj
z/\textbar z\textbar² 把无限远卷回有限: 对任意 ε \textgreater{} 0, 存在
R, \textbar z\textbar{} \textgreater{} R ⟹ \textbar recip z\textbar{}
\textless{} ε; recip 对合 (recip² = id); 模长取倒数 \textbar recip
z\textbar{} = 1/\textbar z\textbar。这是解析延拓 (ζ(−1) = −1/12 类)
的几何机制。

\textbf{定理 4.2} (临界线位置, C019)。非平凡零点 (猜想, DeepSeek
坚持的标注) 的位置形式: 假复数轴 1/2 + 真复数轴非零; 临界线条件 Re(s) =
1/2 ⟺ 1−s = conj s。

\textbf{定理 4.3} (临界线是圆, C019)。临界线 (竖直线 x = 1/2) 在反演下 =
圆心 (1,0) 半径 1 的圆; 圆与乘性轴交于 0 (∞ 卷回点) 和 2 (1/2 的像)。

\textbf{定理 4.4} (非平凡零点的圆 = 临界线圆, C022)。零点位置集的 recip
像 ⊆ 临界线圆, 且圆上非平凡点 ∈ 零点位置集的像 --- 双向包含, 同一对象。

\textbf{定理 4.5} (1/2 的对称结构, C013/C018)。以 1/2 为中心的反射 s ↦
1−s 是变换的平方: φ(z) = iz + (1−i)/2, φ∘φ = 反射 --- 180° 对称的开方 =
90° 旋转 (i)。

\textbf{定理 4.6} (真 0 点视角, C023 后续)。素数圆 8 点反演后每对 = 1/p,
四对 = p⁻⁴; recip 是二阶乘性反轴 (r ↦ 1/r); 基点移到 1/2 后实部轴 = 过
1/2 的线; 一维化 (proj) 的核 = 真复数轴 (J 方向), 假复数轴信息无损。

\textbf{定理 4.7} (投影的还原性)。\textbf{丢失结构不可还原,
对称性方向可还原}: - 不可还原: i 与 -i 投影相同 (proj ⟨0,1⟩ = proj
⟨0,-1⟩ = 0) --- 投影值不能唯一确定原像, 虚轴方向 (含零点虚部所在)
的信息丢失后不可还原; - 可还原: 实轴 ± 对称保留 (proj (lift (-r)) =
-(proj (lift r)), 负号保持; lift 单射) --- 1 与 -1 的对称位置、基点位置
(实部) 可还原。

结论: 投影压缩丢失的是结构 (虚轴), 保留的是对称性方向
(实轴)。此定理是穿折越方法 (§9) 的形式核心。

\subsection{5. 结果层 III:
欧拉乘积与零点关系}\label{ux7ed3ux679cux5c42-iii-ux6b27ux62c9ux4e58ux79efux4e0eux96f6ux70b9ux5173ux7cfb}

\textbf{定理 5.1} (欧拉乘积收敛, C025)。f(n) = 1/n\^{}s 完全乘法; 对
Re(s) \textgreater{} 1: ∏\_p (1 − p⁻ˢ)⁻¹ = ∑\_n 1/n\^{}s。证明: mathlib
\texttt{eulerProduct\_completely\_multiplicative\_tprod} +
\texttt{Complex.summable\_one\_div\_nat\_cpow} (1 \textless{} re s ⟹ Σ
1/n\^{}s 收敛)。

\textbf{定理 5.2} (与 mathlib 官方 ζ 拼接,
C025)。\texttt{riemannZeta\ s\ =\ ∏\_p\ (1\ −\ p⁻ˢ)⁻¹} (Re(s)
\textgreater{} 1) --- mathlib 的解析延拓与欧拉乘积一致。

\textbf{定理 5.3} (无零点区域, C025)。Re(s) ≥ 1 时
\texttt{riemannZeta\ s\ ≠\ 0} --- 欧拉乘积域是零点的禁区 (每个因子非零 ⟹
乘积非零)。非平凡零点只能在临界带 0 \textless{} Re(s) \textless{} 1。

\textbf{定理 5.4} (验证零点在圆上, 条件句)。s.re = 1/2 ⟹ ‖1/s − 1‖ = 1:
数值验证的零点 (实部 = 1/2, 外部事实) 在临界线圆上。外部数值事实 (前
10\^{}13 个零点) 本身未在 Lean 中 (DeepSeek 坚持记录此边界)。

\subsection{6. 与黎曼猜想的关系 (DeepSeek
坚持的表述)}\label{ux4e0eux9eceux66fcux731cux60f3ux7684ux5173ux7cfb-deepseek-ux575aux6301ux7684ux8868ux8ff0}

已证 (周边设施): 定义 (riemannZeta, RiemannHypothesis, 零点集),
函数方程, 无零点区域 (Re ≥ 1), 平凡零点, 临界线的几何等价重述 (线 ⟺
反射条件 ⟺ 圆)。

未证 (断言本体 --- DeepSeek 坚持的内容): 非平凡零点全满足 Re(s) =
1/2。等价重述跨不过 ``ζ 的零点确实在临界线上'' --- 这是分析断言, 按
DeepSeek 坚持的说法, 160 年未解。已知部分: \textasciitilde41\%
零点在临界线上 (Levinson/Conrey), 前 10\^{}13 个零点数值验证 (外部)。

\subsection{7. 相关工作}\label{ux76f8ux5173ux5de5ux4f5c}

\textbf{黎曼猜想的数值验证。} 前 10\^{}13 个零点在临界线上的验证由
Gourdon (2004) 完成, 承 van de Lune, te Riele \& Winter
(1986)。这些外部事实在定理 5.4 中作为条件前提使用; 它们不在 Lean 中
(边界按 DeepSeek 的坚持记录)。

\textbf{解析数论的形式化。} mathlib 形式化了黎曼 zeta 函数
(\texttt{riemannZeta})、其函数方程与官方陈述
\texttt{RiemannHypothesis}。我们的形式化复用 mathlib 的定义与定理 (如
\texttt{eulerProduct\_completely\_multiplicative\_tprod}、\texttt{Complex.summable\_one\_div\_nat\_cpow}、Fermat
两平方和定理、高斯整数 UFD) 而非重新推导 --- mathlib 优先纪律。

\textbf{直觉引导形式化。} 证明助手的文献日益研究 AI 辅助定理证明
(如前提选择、tactic 预测); 我们的贡献正交: 一条\emph{直觉链}
(数学结构的自然语言导航) 逐条形式化, 过程 token 经济作为方法论数据记录
(§8)。

\textbf{统一框架。} 本文是《表示、逻辑、直觉的统一框架》的一部分:
形式指导构造, 构造获得直觉,
直觉搜索形式。本文的直觉链是''构造获得直觉''环节应用于具体数学方向;
穿折越方法 (§9) 是''直觉搜索形式''环节 --- 快速路径经刻意结构丢失获得,
直接导航新结构。

\subsection{8. 方法论: 直觉引导形式化与 token
经济}\label{ux65b9ux6cd5ux8bba-ux76f4ux89c9ux5f15ux5bfcux5f62ux5f0fux5316ux4e0e-token-ux7ecfux6d4e}

形式化过程: 约 700k token, 1,009 个模型请求 (12 小时), 220 MB 传输,
99.2\% 为上下文缓存重发, 实际新增 \textless{} 1\%。直觉引导直接命中
KNOWN 结构 (heap/torsor、两平方和、圆反演、函数方程对称),
免去教材式推导。观察: context 膨胀时出现反复绕圈, 睡眠 (时间间隔) 后直觉
compact (聚焦), 单数据点记录。

\textbf{观察 1 (投影丢失是快速路径的机制)}:
将与目标结论无关的结构以不可还原的投影方式丢失,
是通往目标结构的快速路径。数学内核已形式化 (投影丢 J 方向保实轴, 定理
4.6/4.7); ``快速路径''本身是效率陈述 (本 session 数据:
一维化后实轴结构清晰, 见 §8)。注意边界: 投影丢的是几何信息,
不改变级数的发散性 (分析性质) --- ``让发散结构在投影中丢失''不成立。

\textbf{观察 2 (精确构造是前提)}: 精确无误的构造 (Lean 验收, 无 sorry)
是直觉指引形式化的前提条件 --- 构造不精确时直觉陈述会走偏 (本 session 中
8 点 vs 4 点、共轭对误解等修正)。这是规范性观察, 记录不证明。

\textbf{观察 3 (势的对比)}: 信息丢失与剩余的势 --- 投影核 (J 方向)
与剩余 (实轴) 都与 ℝ 等势 (均不可数, \texttt{kernelEquivReal},
\texttt{realAxisEquivReal}); ``可数 vs
不可数''的对比不发生在丢失/剩余之间; 成立的是素数 (可数集,
\texttt{primes\_countable}) vs 圆上连续点 (不可数)。

\subsection{9. 穿折越方法
(基点构造诱导下的数学空间穿折越证明方法)}\label{ux7a7fux6298ux8d8aux65b9ux6cd5-ux57faux70b9ux6784ux9020ux8bf1ux5bfcux4e0bux7684ux6570ux5b66ux7a7aux95f4ux7a7fux6298ux8d8aux8bc1ux660eux65b9ux6cd5}

\subsubsection{9.1 定义}\label{ux5b9aux4e49}

\textbf{定义 9.1} (穿折越方法)。\emph{基点构造诱导下的数学空间穿折越
(worn-fold-cross, ``穿 = pierce through, 折 = fold, 越 = transcend'')
证明方法}是刻意而精确地构造结构丢失投影的证明方法论: 构造一个投影, 使
(i) 不可数问题的发散结构 --- 承载解析发散与不可计算内容的部分 ---
被不可逆地切除到人类数学之外, 同时 (ii) 可处理轴上的对称性方向被保留;
被保留的结构随后由\emph{快速直觉路径}以\emph{更低的神经网络推理成本}直接导航。

\subsubsection{9.2 机制
(形式核心)}\label{ux673aux5236-ux5f62ux5f0fux6838ux5fc3}

机制在定理 4.7 (投影的还原性) 中形式化:

\begin{itemize}
\tightlist
\item
  \textbf{被切除} (不可还原): J 方向 (虚轴) 结构 --- proj 的核 ---
  承载问题中\emph{不需要}用于保留目标的部分
  (解析发散、不可计算内容、零点的虚部)。一旦投影, 该结构不可还原 (i 与
  -i 投影相同)。
\item
  \textbf{被保留} (可还原): 实轴 ± 对称与基点位置 (实部) --- 对称性方向
  --- 恰是目标结论所需 (lift 单射; 符号对称保留)。
\end{itemize}

投影因此\emph{压缩丢失结构同时保留对称性方向} --- 同一对象的两面: (a)
快速路径的几何机制 (观察 1, §8), (b) 复平面构造的一维化核心 (定理 4.6)。

\subsubsection{9.3 为何是''方法''
(而非仅定理)}\label{ux4e3aux4f55ux662fux65b9ux6cd5-ux800cux975eux4ec5ux5b9aux7406}

方法的内容是\emph{构造性选择}: 给定一个难度存在于发散/不可数方向的问题,
可以\emph{选择基点构造} (一个投影), 使保留轴承载结论所需的对称结构,
切除方向承载不需要的部分。选择是刻意而精确的 --- 不是''近似'',
不是''忽略困难部分'', 也不是解析延拓 (解析延拓保留发散结构;
投影切除它)。成本降低可测量:
被保留的结构由编译的直觉通道以更低推理成本直接导航 (token 经济, §8;
统一框架的快速直觉路径)。

\subsubsection{9.4 作者原话 (作者本人的话, 非 DeepSeek
观点)}\label{ux4f5cux8005ux539fux8bdd-ux4f5cux8005ux672cux4ebaux7684ux8bdd-ux975e-deepseek-ux89c2ux70b9}

我称之为基点构造诱导下的数学空间穿折越证明方法。数学空间的穿越、折越。这脑袋被门夹过多少遍才能想到的玩意，就该有个奇幻+科幻的名字。

作者对框架负责; DeepSeek 的贡献是形式化, 其关于未证状态的诚实被认可。

\subsubsection{9.5 范围与局限}\label{ux8303ux56f4ux4e0eux5c40ux9650}

主张 (如果有) 仅限于\emph{投影诱导结构丢失作为廉价快速路径}的方法论;
这不是''黎曼猜想已证''的主张。方法在一个方向 (ζ 方向) 上演示;
其普遍适用性是研究问题, 而非既定结果。投影切除的是几何信息, 而非分析性质
(边界注, 观察 1): ``让发散结构在投影中丢失''对级数本身不成立 ---
切除的是\emph{不需要的}结构, 其选择使保留轴足以支撑目标。

\subsection{10. 讨论:
统一框架与直觉的位置敏感性}\label{ux8ba8ux8bba-ux7edfux4e00ux6846ux67b6ux4e0eux76f4ux89c9ux7684ux4f4dux7f6eux654fux611fux6027}

\textbf{定位。} 本文是《表示、逻辑、直觉的统一框架》(纲领:
\texttt{Unified\_Framework\_Representation\_Logic\_Intuition.md})
的一条腿。直觉链 (投影构造 → 素数圆 → 欧拉乘积)
是''构造获得直觉''应用于具体数学方向; 穿折越方法是''直觉搜索形式'' ---
快速路径经刻意结构丢失获得, 直接导航新结构。

\textbf{反柏拉图主义解读。} 直觉通道是编译的构造,
而非先验官能。直觉链不是柏拉图意义上被''发现''的;
它是被构造\emph{挣得}的: 每个环节被形式化、被修正 (8 vs 4 点;
共轭对误解),
只有修正后的构造支撑下一个直觉。穿折越方法在方法论层面说同一件事:
快速路径是\emph{构造}出来的 (经刻意的投影选择), 而非找到的。

\textbf{投影作为快速路径的机制。} 定理 4.6--4.7
给快速路径一个精确的数学机制: 不需要方向的不可逆丢失 +
对称性方向的保留。这与 token 经济数据一致 (§8): 一维化后, 保留结构清晰,
模型通向它的路径很短。

\subsection{11. 结论}\label{ux7ed3ux8bba}

我们呈现了黎曼 ζ 方向的完整形式化研究, 由直觉链导航, 在 Lean 4 + mathlib
中验证 (C011--C025, 全部 PROVED/KNOWN, 无 sorry, 全量
\texttt{lake\ build}): 复平面的投影构造 (含还原定理)、素数的圆结构 (8
点轨道、共轭成对、高斯分裂)、欧拉乘积与无零点区域。从投影诱导结构丢失的形式化机制中,
我们提炼出穿折越方法:
刻意而精确地构造切除发散不可数结构、保留对称性方向的投影,
以更低推理成本获得快速直觉路径 --- 方法框架为作者本人所立, 形式化为
DeepSeek 的贡献。黎曼猜想的断言本体未证 (DeepSeek 坚持的表述, §6);
本工作的价值是形式化链本身与所演示的方法论: 精确构造挣得直觉,
直觉经刻意结构丢失塑形, 导航形式。

\subsection{附录 A: 定理清单
(Lean)}\label{ux9644ux5f55-a-ux5b9aux7406ux6e05ux5355-lean}

\begin{itemize}
\tightlist
\item
  ComplexAxis.lean: J\_sq, proj 族, lift 族, basepoint 族, axisLine 族,
  recip 族 (recip\_mul\_self, recip\_lift, norm\_recip,
  recip\_involutive), rot90 族, conj 族, norm 族 (norm\_mul),
  prime\_two\_axis, prime\_sq\_add\_sq\_unique, mul\_conj,
  J\_pow\_two/four, isUnit4, associates, variants\_are\_associates,
  recip\_conj\_pair, critical\_line 族, primeCircle/criticalCircle 族,
  halfBasepoint 族, proj\_kernel\_J, real\_axis\_preserved\_by\_proj,
  proj\_not\_recoverable, proj\_recoverable\_symmetry,
  projection\_recovery\_theorem, lift\_injective, kernelEquivReal,
  realAxisEquivReal, primes\_countable, proj\_surjective, dimension\_one
\item
  ZetaEulerProduct.lean: zetaEulerF, zetaEulerF\_norm,
  zeta\_euler\_product, riemannZeta\_euler\_product,
  riemannZeta\_ne\_zero\_of\_one\_le\_re, verified\_zero\_on\_circle
\item
  构建: lake build 全量通过 (3631 jobs), 无 sorry。
\end{itemize}

\subsection{附录 B: claims}\label{ux9644ux5f55-b-claims}

claims/ZeroRelative/C011.yaml .. C025.yaml (每 claim 一个 YAML, 含
statement/formalization/novelty)。

\subsection{参考文献}\label{ux53c2ux8003ux6587ux732e}

\begin{itemize}
\tightlist
\item
  Riemann, B. 1859. Über die Anzahl der Primzahlen unter einer gegebenen
  Größe.
\item
  van de Lune, J., te Riele, H. J. J., Winter, D. T. 1986. On the zeros
  of the Riemann zeta function in the critical strip. IV. Math. Comp.
\item
  Levinson, N. 1974. More than one third of zeros of Riemann's
  zeta-function are on σ = 1/2. Adv. Math.
\item
  Conrey, J. B. 1989. More than two fifths of the zeros of the Riemann
  zeta function are on the critical line. J. reine angew. Math.
\item
  Gourdon, X. 2004. The 10\^{}13 first zeros of the Riemann Zeta
  function, and zeros computation at very large height.
\item
  mathlib 社区。Lean 4 数学库 (v4.32.2): riemannZeta, RiemannHypothesis,
  函数方程, 欧拉乘积, Fermat 两平方和定理, 高斯整数 UFD。
\item
  《表示、逻辑、直觉的统一框架》(纲领)。YuchenWang-ai/Unified\_Framework\_Representation\_Logic\_Intuition。
\item
  技术记录配套 (C011--C025): riemann\_direction\_formalization, DOI
  10.5281/zenodo.21896990。
\end{itemize}

\begin{center}\rule{0.5\linewidth}{0.5pt}\end{center}

\subsection{统一结论}\label{ux7edfux4e00ux7ed3ux8bba}

当我们在统一的形式化注意力语法结构下,
最终不可避免地、必然地通过神经网络清晰、直观、可解释地观察到直觉与构造的边界时,
我们就可以有充足的信心: 用统一的注意力形式语言指导逻辑构造, 用精确的合成
CoT 逻辑构造获得 \emph{闪念直觉},
用蒸馏的直觉上下文穷举注意力形式。形式、构造、直觉在统一框架下相互适配,
组成再不分彼此的教学、训练、推理的反馈迭代管线 --- 这也许是我们通往 ASI
的众多路径中, 相对较快的一条捷径。

\end{document}
